\chapter*{Resumen}
\phantomsection
    Las lesiones cutáneas malignas plantean un reto diagnóstico, especialmente en contextos donde escasean los datos clínicos etiquetados y los conjuntos públicos proceden de poblaciones distintas a la de Santander. Este trabajo desarrolla un algoritmo de clasificación de lesiones cutáneas mediante aprendizaje autosupervisado y estudia, como líneas complementarias, la segmentación y la transformación y estimación del tono de piel. En la comparación congelada sobre HAM10000, los intervalos bootstrap pareados del 95\,\% de B71 frente a ocho líneas base excluyeron cero; estos contrastes no tuvieron una corrección global por multiplicidad y se interpretan dentro de esa familia. La adaptación con imágenes sin etiquetas de PAD-UFES-20 mostró una mejora piloto frente al control de cómputo, aunque el modelo adaptado permaneció por debajo del codificador público. La variante S4-A2 obtuvo Dice de 0,7221 en ISIC2018 y 0,7949 en HAM10000, sin demostrar una mejora general del clasificador al integrar máscaras. En la evaluación propia de PAD-UFES-20, el estimador ITA quedó por debajo de la línea base por mayoría. Los resultados son experimentales y no constituyen validación clínica ni evidencia de representatividad regional. La implementación del prototipo previsto en el objetivo 4 sigue pendiente mientras se seleccionan el modelo y la configuración finales.
    \vspace{0.5cm}
    \textit{Palabras clave:} aprendizaje autosupervisado, SimSiam, DINOv2, clasificación de lesiones cutáneas, adaptación de dominio, segmentación de lesiones, estimación del tono de piel.

\chapter*{Abstract}
\phantomsection
    Malignant skin lesions pose a diagnostic challenge, particularly in settings with limited labeled clinical data and public datasets collected from populations that differ from Santander. This study develops a self-supervised skin lesion classification algorithm and examines segmentation and skin-tone transformation and estimation as complementary workstreams. In the frozen comparison on HAM10000, the paired bootstrap 95\% confidence intervals for B71 versus eight baselines excluded zero. These contrasts were not subject to a global multiplicity correction and are interpreted within that comparison family. Adaptation using unlabeled PAD-UFES-20 images showed a pilot improvement over the compute-matched control, although the adapted model remained below the public encoder. S4-A2 achieved Dice scores of 0.7221 on ISIC 2018 and 0.7949 on HAM10000, without demonstrating an overall classifier improvement from mask integration. In the project’s PAD-UFES-20 evaluation, the ITA estimator remained below the majority-class baseline. The findings are experimental and do not constitute clinical validation or evidence of regional representativeness. Implementation of the prototype planned under objective 4 remains pending while the final model and configuration are selected.
    \vspace{0.5cm}
    \textit{Keywords:} self-supervised learning, SimSiam, DINOv2, skin lesion classification, domain adaptation, lesion segmentation, skin tone estimation.
